\begin{lemma}[Adjoining a fixed exact edge is unique]
Let $T$ be an ordered tree-sequence in a loopless multigraph, and let $f$
be one fixed edge identity.  If both $U$ and $W$ satisfy
\noderef{TreeSequenceAdjoinEdge}$(T,f,-)$, then $U=W$ as full ordered
edge-valued tree-sequences.
\end{lemma}
\begin{proof}
Write $n=|T|$.  Unpacking \noderef{TreeSequenceAdjoinEdge} for $U$ and
$W$ gives
\[
 |U|=|W|=n+1,
 \qquad T\preccurlyeq U,
 \qquad T\preccurlyeq W,
\]
where both initial-segment relations are the exact ordered relation of
\noderef{TreeSequenceInitialSegment}, and the final edge entry of both
$U$ and $W$ is the same edge identity $f$.

For every vertex index $i\leq n$, the two initial-segment vertex clauses
give $U(i)=T(i)=W(i)$.  Likewise, for every edge index belonging to $T$,
the edge clauses give $U(i)=T(i)=W(i)$.  The remaining final edge entries
also agree because each is exactly $f$.  It remains only to identify the
two new final vertices.

By the attachment clause of \noderef{TreeSequence}, there are indices
$q_U,q_W\leq n$ such that
\[
  \operatorname{ends}(f)=
    \{U(n+1),T(q_U)\}
  \quad\hbox{and}\quad
  \operatorname{ends}(f)=
    \{W(n+1),T(q_W)\}.                                      \tag{1}
\]
Equality of the two unordered endpoint pairs in (1) has two possible
orientations.  In the direct orientation it immediately gives
$U(n+1)=W(n+1)$.  In the swapped orientation it gives
$U(n+1)=T(q_W)$.  But the initial-segment vertex equality also gives
$T(q_W)=U(q_W)$, so this would identify the new vertex entry $U(n+1)$
with an earlier vertex entry of $U$.  This contradicts the vertex
injectivity clause of \noderef{TreeSequence}.  Hence only the direct
orientation is possible, and the final vertices agree.

Thus $U$ and $W$ have the same length, the same vertex entry at every
index, and the same exact edge identity at every index.  All remaining
fields of \noderef{TreeSequence} are propositions asserting positivity,
injectivity, the first-edge condition, and attachment.  Proof irrelevance
identifies those fields, so the complete structures satisfy $U=W$.
The argument compares the identity $f$ itself, not merely its endpoint
pair, and therefore remains valid when parallel edges are present.
\end{proof}
